% =====================================================================
%  Chapter 34 — The Window Sensitivity
% =====================================================================

\chapter{The Window Sensitivity}
\label{ch:window}

Every backtest is a story about a window. This chapter takes the book's
conclusion --- that the governor contributes and the model does not --- and
asks what happens to it when the window moves. It is the check that
separates a finding from a coincidence, and it is the one that a reader
should apply to any backtest they are shown, including this one.

\section{The method}

The method is simple and it is worth stating precisely, because it is not
the same as the usual sensitivity analysis.

The system's layer attribution is recomputed from \numWindowCount{}
different starting months, spread across the record. For each starting
month, the policy alone, the policy with the model, and the policy with
the governor are each re-measured over the remaining months, and the
contribution of each layer is the difference in compound growth rate
between the layer and the policy alone.

The start dates are twelve months apart, and a window is only counted if
it contains at least sixty months of record, because a compound growth
rate computed over three years is not a compound growth rate in any
useful sense.

\section{What the governor does across the windows}

The governor's contribution to the compound growth rate is positive in
every one of the \numWindowCount{} windows. Its mean is \numGovMean{} a
year, its minimum is \numGovMin{}, and its maximum is \numGovMax{}.

That consistency is the finding. It is not a large effect, and the book
does not claim it is. But it is the same sign and a similar magnitude in
every sub-period of a twenty-four-year record, including the windows
that contain the 2008 crisis and the windows that contain only the
post-crisis recovery.

The layer attribution from every starting month is shown in \figref{fig:window-sensitivity}.

The explanation is not hard to see once the mechanism is understood. The
governor acts on the portfolio's own drawdown, and every sub-period of
the record contains a drawdown. Whatever the sub-period, there is
something for the governor to do, and what it does --- reduce exposure
during a decline --- is worth a little in that sub-period because the
sub-period's subsequent recovery is what it protects.

\section{What the model does across the windows}

The model's contribution has a mean of \numTiltMean{} a year and a median
of \numTiltMedian{}, and it is positive in only \numTiltPosFrac{} of the
windows. Its best window is \numTiltMax{} and its worst is
\numTiltMin{}.

A layer whose contribution is positive in a quarter of the windows and
whose best case is a twentieth of a percentage point a year is a layer
that is not doing anything. The consistency of the governor and the
inconsistency of the model are the same phenomenon seen from two sides: a
mechanism that responds to realised state has a predictable effect on
every sub-period, and a mechanism that responds to a forecast does not.

\section{What this says about backtests in general}

The generalisation is worth making, because the reader will meet it
again.

A backtest whose conclusion is robust to the choice of window is
reporting something about the mechanism. A backtest whose conclusion
depends on the window is reporting something about the sample. The
distinction cannot be made from a single equity curve, and no amount of
significance testing on the full sample substitutes for it.

The test is cheap --- recompute the thing from a dozen starting dates ---
and it is not commonly done. The reason it is not commonly done is the
same reason Chapter 4's five leaks are not commonly checked: the failure
is invisible unless someone goes looking for it.

\begin{takeaway}[The general lesson]
Recompute the conclusion from a dozen starting months. If the sign of the
effect survives, you have a mechanism. If it does not, you have a sample,
and the equity curve you are about to believe is a story about which
months you happened to include.
\end{takeaway}

\section{The limitation of this check}

Two limitations, stated so that the check is not mistaken for more than it
is.

It does not address the regime problem. A window sensitivity checks
whether the conclusion is stable across sub-periods of the record; it
does not check whether the record contains the regime the system will
face next. Chapter 36 takes that up.

And it does not address the family. The window sensitivity as computed
here uses the default configuration. A configuration that was chosen
because it worked in one window would pass its own window sensitivity if
the reader did not also move the start date when they chose it --- which
is exactly the researcher-level leak of Chapter 16, and exactly why the
family is the unit of evaluation rather than the single configuration.

\begin{keytakeaways}
\item The governor's contribution is positive in all \numWindowCount{}
windows, with a mean of \numGovMean{} a year and a minimum of
\numGovMin{}.
\item The model's contribution is positive in only \numTiltPosFrac{} of
the windows, with a median of \numTiltMedian{}.
\item A conclusion that survives a dozen starting months is about a
mechanism; a conclusion that does not is about a sample.
\item The check does not address regime change, and it does not protect
against a configuration that was chosen using the window it was tested
on.
\end{keytakeaways}
